\uuid{n4HE}
\titre{Système linéaire à paramètres/Méthode du Pivot de Gauss}
\auteur{hanani ; mzari}
\organisation{exo7}
\licence{CC BY-NC-SA 4.0}
\datecreate{2020-01-20}
\niveau{L1}
\module{Algèbre}
\chapitre{Déterminant, système linéaire}
\sousChapitre{Système linéaire, rang}
\difficulte{3}
\theme{Systèmes d'équations linéaires, Système linéaire à paramètres, Méthode du Pivot de Gauss}
\exo7qcmid{lille-2.systemes.009}

\contenu{
\begin{question}
\qclassification{200.03}
\qtags{motcle=Système linéaire à paramètres/Méthode du Pivot de Gauss}
On considère le système d'équations, d'inconnue $(x,y,z,t)\in \Rr^4$ et de paramètres des réels
$a,b,c$ et $d$ :
$$(\mathtt{S}) \; \left\{\begin{array}{rcc}
x+y&=&a\\
y+z&=&b\\
z+t&=&c\\
t+x&=&d.\end{array}\right.$$
Quelles sont les assertions vraies ?
\begin{answers}
\bad{$(\mathtt{S}) \Leftrightarrow
\left\{\begin{array}{rcc}
x+y&=&a\\
y+z&=&b\\
z+t&=&c.\end{array}\right.$}
\good{$(\mathtt{S})$ admet une  solution si et seulement si $a+c=b+d$.}
\bad{$(\mathtt{S})$ admet une  solution si et seulement si $a+b=c+d$.}
\good{Le rang de $(\mathtt{S})$ est $3$.}
\end{answers}
\begin{explanations}
$$(\mathtt{S}) \Leftrightarrow
\left\{\begin{array}{rcc}
x+y&=&a\\
y+z&=&b\\
z+t&=&c\\
0&=&b+d-a-c.\end{array}\right.$$
On en déduit que si $a+c\neq b+d$, $(\mathtt{S})$ n'admet pas de solution et que si $a+c= b+d$,
$(\mathtt{S})$ en admet une infinité.
\end{explanations}
\end{question}
}
